\section{Постановка задачи встречи}

\subsection{Орбитальные элементы}

Для описания орбитального движения широко используются кеплеровы элементы орбиты:
\begin{equation*}
    \mathbf{y} = \left( a, \; e, \; \omega, \; i, \; \Omega, \; M \right)^\top,
\end{equation*}
где $a$ --- большая полуось, $e$ --- эксцентриситет, $\omega$ --- аргумент перицентра, $i$ --- наклонение, $\Omega$ --- долгота восходящего узла, $M$ --- средняя аномалия. Вместо средней аномалии $M$ также могут использоваться истинная аномалия $\nu$ и эксцентрическая аномалия $E$. 

Кеплеровы элементы имеют наглядную геометрическую интерпретацию и широко используются в небесной механике, каталогах орбитальных данных и прикладных баллистических программных комплексах. Однако при круговых и экваториальных орбитах они становятся вырожденными: при $e=0$ теряет смысл аргумент перицентра $\omega$, а при $i=0$ становится неопределённой долгота восходящего узла $\Omega$. В численных задачах оптимизации такие особенности могут приводить к скачкам элементов, ухудшению обусловленности и снижению устойчивости итерационных методов.

В настоящей работе для описания орбитального движения КА используются модифицированные равноденственные элементы \cite{danielson1995semianalytic}:
\begin{equation*}
    \mathbf{y} = \left( a, \; h, \; k, \; p, \; q, \; \lambda \right)^\top,
\end{equation*}
\begin{equation*}
    h = e \sin (\omega + I \Omega), \quad
    k = e \cos (\omega + I \Omega),
\end{equation*}
\begin{equation*}
    p = \left[ \tan \frac{i}{2} \right]^I \sin \Omega, \quad
    q = \left[ \tan \frac{i}{2} \right]^I \cos \Omega, \quad
    \lambda = M + \omega + I \Omega,
\end{equation*}
где $I$ --- ретроградный фактор. Помимо средней долготы $\lambda$ далее используется истинная
долгота $L = \nu + \omega + I \Omega$. Преимущество модифицированных равноденственных элементов состоит в том, что они не имеют особенностей. Благодаря этому модифицированные равноденственные элементы широко используются в задачах оптимизации орбитального движения и перелётов с малой тягой \cite{betts2000very, graham2015minimum, kechichian2018applied}.

\subsection{Постановка задачи}

Рассмотрим движение КА на интервале $t \in [t_i, t_f]$, где $t_i$ и $t_f$ --- начальный и конечный моменты времени. Уравнение движения может быть записано в виде:
\begin{equation}
    \dot{\mathbf{y}} = \mathbf{f}(t, \mathbf{y}) + B(\mathbf{y}) \mathbf{a}(t),
\label{dot_y_f_B_a}
\end{equation}
где $\mathbf{f}(t, \mathbf{y}) \in \mathbb{R}^{6}$ отвечает за пассивное движение под воздействием произвольного набора сил, $B(\mathbf{y}) \in \mathbb{R}^{6 \times 3}$ --- матрица чувствительности орбитальных параметров по ускорению в орбитальной системе координат (ОСК), $\mathbf{a}(t) \in \mathbb{R}^{3}$ --- управляющее ускорение в ОСК. Выражения для матрицы $B(\mathbf{y})$ представлены в Приложении А.

Краевые условия можно записать в следующем виде:
\begin{equation}
    \mathbf{y}(t_i) = \mathbf{y}_i, \quad \mathbf{y}(t_f) = \mathbf{y}_t,
\label{border_cond}
\end{equation}
где $\mathbf{y}_i \in \mathbb{R}^{6}$ --- начальные орбитальные параметры, $\mathbf{y}_t \in \mathbb{R}^{6}$ --- целевые орбитальные параметры.

Объединив уравнение движения (\ref{dot_y_f_B_a}) и краевые условия (\ref{border_cond}), задачу встречи можно записать в виде
\begin{equation}
\begin{aligned}
    \min_{\mathbf{a} \in \mathcal{A}} \quad &\int_{t_i}^{t_f} \Vert \mathbf{a}(t) \Vert_2 ~ dt, \\
    \text{s.t.} \quad &\dot{\mathbf{y}} = \mathbf{f}(t, \mathbf{y}) + B(\mathbf{y}) \mathbf{a}(t), \\
    &\mathbf{y}(t_i) = \mathbf{y}_i, \\
    &\mathbf{y}(t_f) = \mathbf{y}_t,
\end{aligned}
\label{rend_genaral_form}
\end{equation}
где $\mathcal{A}$ --- множество допустимых управлений, задаваемое ограничениями на величину, направление, длительность и допустимые интервалы маневрирования.

\subsection{Пассивные силы}

Функция пассивного движения $\mathbf{f}(t, \mathbf{y})$ может быть представлена в виде: 
\begin{equation}
    \mathbf{f}(t, \mathbf{y}) = \mathbf{f}_0 (\mathbf{y}) + B(\mathbf{y}) \mathbf{a}_{\text{pert}}(t, \mathbf{y}),
\end{equation}
где $\mathbf{f}_0 (\mathbf{y}) = \left( \mathbf{0}_{1 \times 5}, n \right)^\top$ --- вклад тяготения центрального поля, $n = \sqrt{\mu / a^3}$ --- среднее движение, $\mathbf{a}_{\text{pert}}(t, \mathbf{y})$ --- ускорение от пассивных возмущающих сил в ОСК. Такое ускорение определяется по второму закону Ньютона:
\begin{equation}
    \mathbf{a}_{\text{pert}} = \frac{\mathbf{F}_{\text{pert}}}{m},
\label{a_F_m}
\end{equation}
где $m$ --- масса КА, а $\mathbf{F}_{\text{pert}}$ --- результирующая пассивных возмущающих сил в ОСК. Представленные далее алгоритмы маневрирования не будут зависеть от конкретной баллистической модели. Но для методической полноты работы и возможности воспроизведения результатов, приведём описание основных возмущающих сил. Подробное обозрение может быть найдено, например, в \cite{vallado2001fundamentals} или \cite{авдюшев2015численное}. Оценку величин различных сил для орбит различной высоты можно найти в \cite{марков2016анализ}. 

Результирующая $\mathbf{F}_{\text{pert}}$ может быть представлена в виде суммы слагаемых:
\begin{equation}
    \mathbf{F}_{\text{pert}} = \mathbf{F}_{\text{grav}} + \mathbf{F}_{\text{drag}} + \mathbf{F}_{\text{other}} + \mathbf{F}_{\text{srp}},
\label{F_passive}
\end{equation}
где $\mathbf{F}_{\text{grav}}$ --- возмущающая составляющая гравитационного притяжения Земли, $\mathbf{F}_{\text{drag}}$ --- сила сопротивления атмосферы, $\mathbf{F}_{\text{other}}$ --- сила притяжения третьих тел, $\mathbf{F}_{\text{srp}}$ --- сила давления солнечного излучения. Обычно силы в (\ref{F_passive}) рассчитываются в геоцентрической инерциальной СК. Для использования формулы (\ref{a_F_m}) их необходимо преобразовать в ОСК. Рассмотрим их вычисление более подробно.

\subsubsection{Гравитационное притяжение Земли}
Земля является телом очень сложной формы с постоянно изменяющейся внутренней структурой. Поэтому использование приближения точечного потенциала в задачах, требующих высокую точность прогноза, может быть некорректным. Наиболее часто используемым подходом является разложение гравитационного потенциала по сферическим функциям, согласно конвенции IERS 2010 \cite{petit20122010}:
\begin{equation}
    U(r, \phi, \lambda_g) = U_{\text{c}}(r) + U_{\text{nc}}(r, \phi, \lambda_g), \quad \quad
    U_{\text{c}}(r) = \frac{\mu}{r},
\label{grav_potential_1}
\end{equation}
\begin{equation}
    U_{\text{nc}}(r, \phi, \lambda_g) = \frac{\mu}{r} \sum_{n=1}^{\infty} \left(\frac{R}{r}\right)^{n} \sum_{m=0}^{n} \bar{P}_{nm} \left( \sin{\phi} \right) \left(\bar{C}_{nm}\cos{(m\lambda_g)} + \bar{S}_{nm} \sin{(m\lambda_g)} \right),
\label{grav_potential_2}
\end{equation}
где $(r, \phi, \lambda_g)$ --- радиус, широта и долгота --- координаты точки в геоцентрической СК, $\mu$ --- гравитационный параметр Земли в соответствии моделью EGM 2008 \cite{pavlis2012development}, $R$ --- экваториальный радиус планеты. $\bar P_{nm}$ --- нормированные присоединенные полиномы Лежандра, $\bar{C}_{nm}$ и $\bar{S}_{nm}$ --- коэффициенты разложения.

Сила $\mathbf{F}_{\text{grav}}$, действующая на тело массы $m$ определяется следующим образом:
\begin{equation*}
    \mathbf{F}_{\text{grav}} = m \nabla U_{\text{nc}}.
\end{equation*}
На практике бесконечный ряд (\ref{grav_potential_1}, \ref{grav_potential_2}) обрезается до некоторого количества членов (гармоник) $N$, которое необходимо для достижения требуемого уровня точности.

\subsubsection{Сопротивление атмосферы}
Сила атмосферного сопротивления является одним из основных возмущающих факторов для КА на низких околоземных орбитах. Её вычисление включает два этапа: определение параметров атмосферы в точке расположения КА и расчёт силы по выбранной модели взаимодействия набегающего потока с поверхностью аппарата.

Модели атмосферы можно условно разделить на статические и динамические. Динамические модели учитывают зависимость плотности от солнечной и геомагнитной активности. К числу распространённых моделей относятся зарубежные JB2006 \cite{bowman2008jb2006}, JB2008 \cite{bowman2008new}, NASA NRLMSISE-00 \cite{picone2002nrlmsise} и отечественная ГОСТ Р 25645.166-2004 \cite{gost25645_166_2004}. В настоящей работе будет использоваться последняя из перечисленных моделей.

Сила сопротивления задаётся выражением
\begin{equation}
    \mathbf{F}_{\text{drag}} = - \frac{1}{2} \rho C_D S_{\text{drag}} \Vert \mathbf{v}_{\text{rel}} \Vert \mathbf{v}_{\text{rel}}, 
    \label{drag_force}
\end{equation}
где $\rho$ --- плотность атмосферы, $C_D$ --- коэффициент сопротивления, $S_{\text{drag}}$ --- площадь миделева сечения, $\mathbf{v}_{\text{rel}}$ --- скорость КА относительно набегающего газового потока.

\subsubsection{Притяжение прочих небесных тел}
Для расчёта силы тяготения от удалённых небесных тел нет необходимости применять сложные модели для аппроксимации их гравитационного поля. Можно получить приемлемую точность расчёта, если положить, что КА взаимодействует с ними как с точечными телами. В таком случае искомая сила  может быть записана в виде следующей суммы:
\begin{equation}
    \mathbf{F}_{\text{other}} = m \sum_{k=1}^{K} \mu_k \left( \frac{\mathbf r_k-\mathbf r}{\Vert \mathbf r_k-\mathbf r \Vert^3} - \frac{\mathbf r_k}{\Vert \mathbf r_k \Vert^3}
    \right),
\label{third_body_gravity}
\end{equation}
где $m$ --- масса КА, $\mathbf{r}$ --- радиус-вектор КА относительно центра Земли, $\mathbf{r}_k$ --- радиус-вектор $k$-го небесного тела относительно центра Земли, $\mu_k$ --- гравитационный параметр этого тела. Положения третьих тел могут быть получены из эфемерид, которые предоставляются JPL NASA: JPL DE430, JPL DE431 \cite{folkner2014planetary}.

\subsubsection{Давление солнечного излучения}

Сила давления солнечного излучения зависит от положения КА относительно Солнца, площади освещённой поверхности и оптических свойств аппарата. Для низких околоземных орбит этот эффект обычно меньше атмосферного сопротивления, однако может становиться существенным при увеличении высоты орбиты, большой площади аппарата или длительном интервале прогноза.

В геоцентрической инерциальной СК силу давления солнечного излучения можно записать в виде
\begin{equation}
    \mathbf F_{\mathrm{srp}}
    =
    \Phi
    P_\odot C_R S_{\mathrm{srp}}
    \left(
        \frac{a_E}{\|\mathbf r-\mathbf r_\odot\|}
    \right)^2
    \frac{\mathbf r-\mathbf r_\odot}
    {\|\mathbf r-\mathbf r_\odot\|},
\label{srp_force}
\end{equation}
где $\Phi \in [0,1]$ --- функция тени, $P_\odot$ --- давление солнечного
излучения на расстоянии одной астрономической единицы, $C_R$ --- коэффициент
отражения, $S_{\mathrm{srp}}$ --- эффективная площадь КА для давления
солнечного излучения, $a_E$ --- астрономическая единица, $\mathbf r$ ---
радиус-вектор КА относительно Земли, $\mathbf r_\odot$ --- радиус-вектор
Солнца относительно Земли. Вектор $\mathbf r-\mathbf r_\odot$ направлен от
Солнца к КА, поэтому сила \eqref{srp_force} направлена от Солнца.

Функция $\Phi$ учитывает возможное затенение КА Землёй или другими небесными
телами. В простейшем случае она принимает значения $0$ и $1$, соответствующие
полному затенению и полной освещённости; более продвинутые модели могут
использовать промежуточные значения при прохождении полутени.

